\documentclass[11pt]{article}

\usepackage[a4paper,margin=27mm]{geometry}
\usepackage{amsmath,amssymb,amsthm,mathtools}
\usepackage[hidelinks]{hyperref}
\usepackage{microtype}
\usepackage{enumitem}

\newtheorem{definition}{Definition}[section]
\newtheorem{proposition}[definition]{Proposition}
\newtheorem{theorem}[definition]{Theorem}
\newtheorem{corollary}[definition]{Corollary}
\newtheorem{example}[definition]{Example}
\newtheorem{remark}[definition]{Remark}

\newcommand{\X}{\mathcal X}
\newcommand{\Acal}{\mathcal A}
\newcommand{\Bcal}{\mathcal B}

\hypersetup{
  pdftitle={WR-III: Geometry of Admissible World Fibres},
  pdfauthor={A. A. Malachevsky},
  pdfsubject={World Realizability and Epistemic Horizons (WREH), WR-III},
  pdfkeywords={world fibres, inverse images, identifiability, rigidity, refinement, set-valued maps, realizability walls}
}

\title{\textbf{WR-III Mathematical Spine v0.2}\\[4pt]
\large Geometry of Admissible World Fibres:\\
Refinement, Rigidity and Realizability Walls}
\author{A. A. Malachevsky\\
ORCID: 0009-0008-6009-3196\\
World Realizability \& Epistemic Horizons (WREH)\\
Research Programme \& Community}
\date{8 October 2026}

\begin{document}
\maketitle

\begin{abstract}
WR-III studies the geometry of nonempty inverse fibres left after the
identifiability and realizability questions of WR-I and WR-II. For a declared
geometric realization carrier and a response map indexed by a protocol bundle,
we introduce the inverse-fibre correspondence over the realizable response
image. We prove monotone fibre shrinkage under protocol refinement, establish
upper hemicontinuity on compact carriers, characterize lower hemicontinuity by
openness of the response map, and derive upper semicontinuity of fibre diameter
on compact metric carriers. These results yield open finite-resolution
rigidity loci. In the smooth proper setting, Ehresmann's theorem supplies local
triviality over regular values, motivating an intrinsic structural wall as the
failure locus of local fibre triviality. This wall is kept distinct from the
boundary of the attainable response image. The framework is model-relative:
no response-space wall is identified with a physical boundary or phase
transition without an additional bridge.
\end{abstract}

\section{Upstream interface}

WR-III begins after WR-I and WR-II. Let \(\X\) be the declared geometric
realization carrier and let
\[
R_{\Acal}:\X\to Y_{\Acal}
\]
be the response map associated with an admissible protocol bundle \(\Acal\).
Its realizable response image is
\[
E_{\Acal}:=R_{\Acal}(\X).
\]

The carrier \(\X\) may be the original admissible world class or a justified
gauge-reduced carrier. WR-III does not assume that a gauge quotient exists, is
Hausdorff, metric, or smooth unless those properties are stated.

\begin{definition}[Admissible world fibre]
For \(e\in E_{\Acal}\), define
\[
\Phi_{\Acal}(e):=R_{\Acal}^{-1}(e).
\]
\end{definition}

Every fibre considered in WR-III is nonempty by construction. Whether a
response is realizable belongs upstream to WR-II.

\section{Refinement geometry}

Suppose \(\Acal\subseteq\Bcal\) and
\[
R_{\Acal}
=
\rho_{\Bcal\Acal}\circ R_{\Bcal}.
\]

\begin{proposition}[Refinement inclusion]
For every \(e_{\Bcal}\in E_{\Bcal}\),
\[
\Phi_{\Bcal}(e_{\Bcal})
\subseteq
\Phi_{\Acal}\bigl(\rho_{\Bcal\Acal}(e_{\Bcal})\bigr).
\]
\end{proposition}

\begin{proof}
If \(x\in\Phi_{\Bcal}(e_{\Bcal})\), then
\[
R_{\Bcal}(x)=e_{\Bcal}.
\]
Applying \(\rho_{\Bcal\Acal}\) gives
\[
R_{\Acal}(x)=\rho_{\Bcal\Acal}(e_{\Bcal}),
\]
so \(x\) lies in the coarser fibre.
\end{proof}

Assume now that \((\X,d)\) is metric.

\begin{definition}[Fibre diameter]
For nonempty \(F\subseteq\X\), define
\[
\operatorname{diam}F
=
\sup\{d(x,x'):x,x'\in F\}.
\]
Set
\[
D_{\Acal}(e)
:=
\operatorname{diam}\Phi_{\Acal}(e).
\]
\end{definition}

\begin{corollary}[Diameter monotonicity under refinement]
For \(\Acal\subseteq\Bcal\),
\[
D_{\Bcal}(e_{\Bcal})
\le
D_{\Acal}\bigl(\rho_{\Bcal\Acal}(e_{\Bcal})\bigr).
\]
\end{corollary}

\begin{remark}
This monotonicity is informational. It does not imply a physical time arrow,
dynamical contraction, entropy law, or renormalization flow.
\end{remark}

\section{Rigidity}

\begin{definition}[Exact and finite-resolution rigidity]
A response \(e\in E_{\Acal}\) is exactly rigid if
\[
\Phi_{\Acal}(e)
\]
is a singleton.

For \(\varepsilon>0\), it is strictly \(\varepsilon\)-rigid if
\[
D_{\Acal}(e)<\varepsilon.
\]
\end{definition}

\begin{proposition}[Refinement preserves rigidity]
If \(e_{\Bcal}\) refines \(e_{\Acal}\), then exact rigidity or strict
\(\varepsilon\)-rigidity of the coarser fibre implies the same property for
the finer fibre.
\end{proposition}

\begin{proof}
This follows from nonempty fibre inclusion.
\end{proof}

\section{Upper hemicontinuity}

A correspondence
\[
\Phi:E\rightrightarrows\X
\]
is upper hemicontinuous at \(e\) if every open \(U\subseteq\X\) with
\[
\Phi(e)\subseteq U
\]
contains all nearby fibres.

\begin{theorem}[Compact-carrier upper hemicontinuity]
\label{thm:uhc}
Let \(\X\) be compact, let \(E\) be Hausdorff, and let
\[
R:\X\to E
\]
be continuous and surjective. Then the inverse correspondence
\[
\Phi(e)=R^{-1}(e)
\]
has nonempty compact values and is upper hemicontinuous.
\end{theorem}

\begin{proof}
Every fibre is nonempty and closed, hence compact. Suppose upper
hemicontinuity fails at \(e\). Then some open \(U\supseteq\Phi(e)\) admits a
net \(e_i\to e\) and points
\[
x_i\in\Phi(e_i)\setminus U.
\]
Compactness gives a convergent subnet \(x_i\to x\in\X\setminus U\). By
continuity,
\[
R(x)=e,
\]
contradicting \(x\in\Phi(e)\subseteq U\).
\end{proof}

\section{Lower hemicontinuity and openness}

A correspondence \(\Phi:E\rightrightarrows\X\) is lower hemicontinuous if for
every open \(U\subseteq\X\), the set
\[
\{e\in E:\Phi(e)\cap U\neq\varnothing\}
\]
is open.

\begin{theorem}[Openness criterion]
Let \(R:\X\to E\) be a surjection. Its inverse correspondence
\[
\Phi(e)=R^{-1}(e)
\]
is lower hemicontinuous if and only if \(R\) is an open map.
\end{theorem}

\begin{proof}
For every open \(U\subseteq\X\),
\[
R(U)
=
\{e\in E:\Phi(e)\cap U\neq\varnothing\}.
\]
Thus the defining condition for lower hemicontinuity is exactly openness of
\(R\).
\end{proof}

\begin{corollary}
If \(\X\) is compact, \(E\) is Hausdorff, and \(R\) is continuous, surjective,
and open, then its inverse correspondence is both upper and lower
hemicontinuous.
\end{corollary}

\section{Fibre-diameter stability}

\begin{theorem}[Diameter upper semicontinuity]
\label{thm:diam-usc}
Let \((\X,d)\) be compact metric, let \(E\) be metric, and let
\[
R:\X\to E
\]
be continuous and surjective. Then
\[
D(e):=\operatorname{diam}R^{-1}(e)
\]
is upper semicontinuous.
\end{theorem}

\begin{proof}
Let \(e_n\to e\). Choose
\[
x_n,x_n'\in R^{-1}(e_n)
\]
with
\[
d(x_n,x_n')
\ge D(e_n)-\frac1n.
\]
Pass to a subsequence realizing the limsup. Compactness gives
\[
x_n\to x,\qquad x_n'\to x'.
\]
Continuity gives \(R(x)=R(x')=e\), and therefore
\[
\limsup_{n\to\infty}D(e_n)
\le d(x,x')
\le D(e).
\]
\end{proof}

\begin{corollary}[Open finite-resolution rigidity loci]
For every \(\varepsilon>0\),
\[
\mathcal R_\varepsilon
=
\{e\in E:D(e)<\varepsilon\}
\]
is open.
\end{corollary}

\begin{corollary}[Exact rigidity is a \(G_\delta\) locus]
The exact rigidity locus
\[
\mathcal R_0
=
\{e\in E:D(e)=0\}
\]
satisfies
\[
\mathcal R_0
=
\bigcap_{n=1}^{\infty}\mathcal R_{1/n}.
\]
\end{corollary}

\begin{remark}
Exact rigidity need not be open. Unique identification can occur at a
structural singularity.
\end{remark}

\section{Structural regularity and walls}

The word wall is used here in a strictly mathematical, model-relative sense.

\begin{definition}[Structural regular response]
A response \(e\in E\) is structurally regular if there exists an open
neighborhood \(U\subseteq E\) of \(e\) such that
\[
R^{-1}(U)\to U
\]
is a locally trivial topological fibre bundle.
\end{definition}

Let \(E_{\mathrm{reg}}^{\mathrm{str}}\) be the union of all structurally regular
neighborhoods.

\begin{definition}[Structural wall]
The structural wall is
\[
\Sigma_{\mathrm{str}}
=
E\setminus E_{\mathrm{reg}}^{\mathrm{str}}.
\]
\end{definition}

\begin{definition}[Attainable-response boundary]
If \(E\subseteq Y\) is embedded in a declared ambient response space \(Y\),
define
\[
\partial_YE
\]
to be its boundary in \(Y\).
\end{definition}

\begin{definition}[Realizability wall]
Relative to the declared ambient response space \(Y\), define
\[
\Sigma_{\mathrm{real}}
=
\Sigma_{\mathrm{str}}\cup\partial_YE.
\]
\end{definition}

\begin{remark}
The two pieces have different meanings. The boundary separates attainable
from unattainable responses in the declared model. The internal structural
wall records failure of local fibre triviality among attainable responses.
Neither is automatically a physical boundary.
\end{remark}

\section{Smooth proper response maps}

Assume \(\X\) and \(Y\) are smooth manifolds without boundary and
\[
R:\X\to Y
\]
is smooth and proper. Let \(E=R(\X)\).

\begin{theorem}[Regular values are structurally regular]
If \(e\in E\) is a regular value of \(R\), then \(e\) is structurally regular.
Consequently,
\[
\Sigma_{\mathrm{str}}
\subseteq
\operatorname{CritVal}(R)\cap E.
\]
\end{theorem}

\begin{proof}
The critical set is closed. Since a proper map between manifolds is closed,
its critical-value set is closed.

Choose \(x\in R^{-1}(e)\). Because \(e\) is a regular value, the
differential \(dR_x\) is surjective. The submersion theorem therefore gives
an open neighborhood \(U_0\subseteq Y\) of \(e\) contained in the image
\(E\). Since the critical-value set is closed and does not contain \(e\),
there is an open neighborhood \(U_1\subseteq Y\) of \(e\) containing no
critical values. Set \(U=U_0\cap U_1\). Then
\[
R:R^{-1}(U)\to U
\]
is surjective, proper, and a submersion. Ehresmann's fibration theorem implies
local triviality.
\end{proof}

\begin{corollary}[Constant smooth fibre type off the wall]
On each connected component of the regular-value region, all fibres are
diffeomorphic.
\end{corollary}

\begin{remark}
Critical values only contain the possible structural wall under these
hypotheses. A critical value need not change every fibre invariant of interest.
\end{remark}

\section{Minimal example}

\begin{example}[Height response on the sphere]
Let
\[
\X=S^2\subset\mathbb R^3,
\qquad
R(x,y,z)=z.
\]
Then
\[
E=[-1,1].
\]
For \(-1<e<1\), the fibre is a circle. At \(e=\pm1\), the fibre is a
singleton.

With the ambient Euclidean metric,
\[
D(e)=2\sqrt{1-e^2}.
\]
Thus exact rigidity occurs precisely at the two structural wall points.
\end{example}

\begin{remark}[Rigidity is not regularity]
The example shows that exact rigidity and structural regularity are logically
distinct. A singleton fibre can occur where the local fibre type becomes
singular.
\end{remark}

\section{Prior-art boundary}

The surrounding mathematical theories are mature:

\begin{itemize}
\item set-valued analysis studies upper and lower hemicontinuity and compact
correspondences;
\item fibre-bundle theory and Ehresmann's theorem control proper submersions;
\item Sard theory and singularity theory analyze critical values;
\item inverse problems and identifiability theory study parameter fibres,
nonuniqueness, and stability.
\end{itemize}

WR-III does not claim invention of these ingredients. Its narrower target is
their organization downstream of the WREH response and realizability
architecture: nonempty admissible world fibres, refinement shrinkage, metric
rigidity, and model-relative response walls.

\section{Open obligations}

Before WR-III can become a preprint:

\begin{enumerate}
\item perform a detailed prior-art audit against set-valued inverse maps,
identifiability geometry, singular inverse problems, and stratified maps;
\item determine useful fibre invariants beyond diameter;
\item study quantitative distance to the structural wall;
\item characterize refinement-induced wall creation, destruction, and
resolution;
\item add at least one nontrivial inverse-problem application;
\item test whether proper stratified maps give the right extension beyond the
smooth Ehresmann layer;
\item build an interactive demo separating rigidity, ambiguity, and wall
crossing.
\end{enumerate}

\section*{Conclusion}

WR-I identifies response-equivalence classes. WR-II determines whether
compatible responses admit global realizations. WR-III studies the geometry
remaining inside those nonempty inverse fibres.

The first exact chain is
\[
\text{refinement}
\Longrightarrow
\text{fibre inclusion}
\Longrightarrow
\text{nonincreasing diameter}.
\]
Compactness controls upper variation of fibres and fibre diameter. In the
smooth proper setting, regular values carry locally trivial fibres, leaving
critical values and attainable-response boundaries as natural containers for
model-relative realizability walls.

No wall introduced here is, without an additional physical bridge, a wall of
spacetime or of the Universe.

\begin{thebibliography}{99}

\bibitem{MalachevskyWRI}
A. A. Malachevsky,
\emph{Response-Defined World Spaces: Response Equivalence,
Finite-Resolution Separation, and Identifiability of Global Realizations},
WR-I, WREH (2026).
DOI: 10.5281/zenodo.23210258.

\bibitem{MalachevskyWRII}
A. A. Malachevsky,
\emph{Finite Consistency and Global World Realizability},
WR-II, WREH (2026).
DOI: 10.5281/zenodo.23224154.

\bibitem{AubinFrankowska}
J.-P. Aubin and H. Frankowska,
\emph{Set-Valued Analysis},
Birkh\"auser, 1990.

\bibitem{Ehresmann}
C. Ehresmann,
\emph{Les connexions infinit\'esimales dans un espace fibr\'e
diff\'erentiable},
Colloque de Topologie, Bruxelles, 1950.

\end{thebibliography}

\end{document}
